\begin{frame}{What is Undefined Behaviour}
  Undefined Behaviour (UB) is a set of \textbf{erroneous} operations that may cause the system to behave badly. LLVM distinguish two categories of UB:
\begin{columns}[T]
    \begin{column}{0.44\textwidth}
        \begin{center}
        \large{\textbf{Immediate UB}}\\
        \end{center}
        Errors that have serious consequences.
        \begin{itemize}
            \item A division by zero
            \item Dereferencing a \texttt{null} pointer
        \end{itemize}
    \end{column}
    \begin{column}{0.44\textwidth}
        \begin{center}
        \large{\textbf{Deferred UB}}\\
        \end{center}
        Operations that may produce unpredictable values.
        \begin{itemize}
            \item Signed integers overflow
            \item The value of an unassigned variable
        \end{itemize}
    \end{column}
\end{columns}
\pdfpcnote{
Introduce the two-category taxonomy. Immediate UB halts the system immediately —
the compiler cannot hoist code past it. Deferred UB produces a bad value but
execution continues — this is what enables speculative optimizations. Keep this
brief; the next slides go into each case in detail.
}
\end{frame}
